\documentclass{beamer}

\usepackage[utf8]{inputenc}
\input{../helper}

\title{Introdução à Computação \\ listas encadeadas}
\author{Alexandre Rademaker}
\date{}

\begin{document}

\begin{frame}
 \maketitle
\end{frame}

\begin{frame}[allowframebreaks,fragile]{listas encadeadas}

  Com as listas encadeadas (linked lists) podemos armazenar valores em
  diferentes posições de memória (não contínuas).

  Com isso, podemos acrescentar novos valores na lista continuamente,
  sem necessidade de estabelecer um limite `a priori' nem realocação
  quando o limite é atingido.

  \framebreak

  Para cada valor, alocamos também o endereço do próximo valor na
  lista, o par é uma \href{https://en.wikipedia.org/wiki/Cons}{cons
    cell}.

  \begin{center}
    \includegraphics[width=.9\textwidth]{linked-list.png}
  \end{center}
  
  \framebreak

  Em C, podemos implementar uma \textbf{cons cell} (node) com o código.
  
  \vspace{.4cm}

  \begin{lstlisting}[style=normalc]
    typedef struct node
    {
      int number;
      struct node *next;
    } t_node;
  \end{lstlisting}

  Trata-se de uma estrutura de dados recursiva.

  \framebreak

  \begin{lstlisting}[style=normalc]
    // lista vazia
    t_node *list = NULL; 

    t_node *n = malloc(sizeof(t_node));
    if (n != NULL) {
      n->number = 5;  // same as (*n).number
      n->next = NULL;
    }
    
    list = n;
  \end{lstlisting}

  A partir da lista vazia, adicionamos o número 5 na frente da lista.

  \framebreak

  Graficamente, podemos representar o resultado do código anterior com

  \begin{center}
    \includegraphics[width=.5\textwidth]{list_with_one_node_n.png}
  \end{center}

\end{frame}


\begin{frame}{lista encadeada na memória}

  \foreach \n in {1,...,25}{
    \includegraphics<\n>[width=.9\textwidth]{ll/ll-\n.png}}

\end{frame}


\begin{frame}{como \textbf{não} fazer}

  \foreach \n in {1,...,5}{
    \includegraphics<\n>[width=.9\textwidth]{ll/error-\n.png}}

\end{frame}

\begin{frame}{como \textbf{não} fazer}

  Não podemos atualizar o ponteiro \ilcode{list} antes de armazenar
  seu valor em \ilcode{n->next}.
  
\end{frame}


\begin{frame}{adicionando elementos no início}

  adicionar um elemento no início tem custo $O(1)$

  \foreach \n in {1,...,6}{
    \includegraphics<\n>[width=.9\textwidth]{ll/ll0-\n.png}}
  
\end{frame}

\begin{frame}{adicionando elementos em uma posição qualquer}

  após localizar o elemento, eventualmente $O(n)$, custo de $O(1)$

  \foreach \n in {1,...,3}{
    \includegraphics<\n>[width=.9\textwidth]{ll/llx-\n.png}}
  
\end{frame}


\begin{frame}[allowframebreaks]{Sumário}

  \includegraphics[width=.9\textwidth]{linked-lists.png}

  \textbf{arrays} também são chamados \textbf{vectors}.

  \framebreak

  complexidade deve ser avaliada para cada operação

  \begin{itemize}
  \item para recuperar o elemento da posição $n$? \\ \textbf{R}: array $O(1)$, lista $O(n)$    
  \item inserir elemento no início \\ \textbf{R}: array $O(n)$, lista $O(1)$
  \item inserir elemento no final \\ \textbf{R}: array $O(1)$/$O(n)$, lista
    $O(n)$
  \item inserir elemento no meio da estrutura?
  \item busca de um elemento na estrutura não ordenada?
  \item busca de um elemento na estrutura ordenada?
  \end{itemize}
  
\end{frame}



\end{document}

%%% Local Variables:
%%% mode: latex
%%% TeX-master: t
%%% End:
